\sporkSection{Trace Introduction}{sec:TraceIntro}

Before introducing the full language,
we introduce the basics of what constitutes a trace,
and describe the reasoning behind our notion of sequential-parallel equivalence.
This will be presented without proof in natural language,
and forms the basis of our language design and synchronization checking.

Informally, the sequence of trace events is ordered in two ways:
in \emph{trace total order}, which is given by the trace sequence order
(i.e. the order a sequential execution would issue the trace events),
and in \emph{parallel partial order},
which orders pairs of trace events only when guaranteed by executed synchronization or by sequencing of non-async instructions within a single thread.

We do not require perfect agreement of the two orderings.
We require only enough to guarantee that each read of an array element $\sporkx[\sporkz^*]$,
as executed by the parallel hardware,
will observe the value stored by the most recent prior mutate of $\sporkx[\sporkz^*]$,
with \emph{recent} and \emph{prior} defined by trace total order (Figure~\ref{fig:RequiredProperty}).

\begin{figure}[h]
\sporkForceStacktrue
\sporkCodeRule
\centering
\begin{sporkfigure}
  % Body of the Required Property figure (26.tex frame RequiredProperty, paper RequiredProperty.tex):
% raw TikZ, \input inside a sporkfigure. Six trace events in trace total order (left to right);
% reveals: blue read-from arrows for x[0] then x[1], then the red dotted "must not" arrows.
% The \sporkSeqReqProp highlights are slide-only. Node names e0..e5 are its interface.
\tikzset{seq event/.style={rectangle, draw, thick, align=center, inner sep=1mm,
  minimum width=15mm}}
\node (e0) [seq event] {\sporkMutateBullet\\\sporkSeqReqProp{0}{0}};
\node (e1) [seq event, right=5mm of e0] {\sporkMutateBullet\\\sporkSeqReqProp{4}{1}};
\node (e2) [seq event, right=5mm of e1] {\sporkMutateBullet\\\sporkSeqReqProp{3}{0}};
\node (e3) [seq event, right=5mm of e2] {\sporkReadBullet\\\sporkSeqReqProp{4}{1}};
\node (e4) [seq event, right=5mm of e3] {\sporkReadBullet\\\sporkSeqReqProp{3}{0}};
\node (e5) [seq event, right=5mm of e4] {\sporkMutateBullet\\\sporkSeqReqProp{0}{0}};

% Square arrows to keep the figure short: blue read-from arrows 3mm out, red ones 6mm out
% above (outside the blue) and 3mm below (no horizontal overlap with the blue x[1]).
\draw [spork reveal=<3->, spork arrow, draw=sporkBlueFg]
  ([xshift=2mm]e2.north) -- ++(0, 3mm) -| ([xshift=2mm]e4.north);
\draw [spork reveal=<4->, spork arrow, draw=sporkBlueFg]
  (e1.south) -- ++(0, -3mm) -| (e3.south);
\draw [spork reveal=<5->, spork arrow, draw=sporkRedFg, densely dotted]
  (e0.north) -- ++(0, 6mm) -| node (olderMid) [pos=0.25, inner sep=0pt] {} ([xshift=-2mm]e4.north);
\node [spork reveal=<5->, above=0.5mm of olderMid] {Must not read \sporkKeyA{older} values};
\draw [spork reveal=<6->, spork arrow, draw=sporkRedFg, densely dotted]
  (e5.south) -- ++(0, -3mm) -| node (futureMid) [pos=0.25, inner sep=0pt] {} (e4.south);
\node [spork reveal=<6->, below=0.5mm of futureMid] {Must not read \sporkKeyA{future} values};

\end{sporkfigure}
\caption{Required property: each read of $\sporkx[\sporkz^*]$ observes the most recent prior mutate of $\sporkx[\sporkz^*]$, with trace total order running left to right.}
\label{fig:RequiredProperty}
\sporkCodeRule
\end{figure}


The requirements we need correspond to resolving the classic hazards:
write-after-write (WAW), read-after-write (RAW), and write-after-read (WAR).
We require agreement between the two orders for pairs of mutates of the same array element $\sporkx[\sporkz^*]$ (WAW),
so that each mutate cleanly overwrites the previous,
and we require order agreement between pairs of read events and mutate events of $\sporkx[\sporkz^*]$,
so that parallel partial order ``squeezes'' a read event between back-to-back mutates (Figure~\ref{fig:RequiredOrderings}).
The executed read is left with no choice but to read one possible value:
the value that would be read in a totally-ordered execution (Figure~\ref{fig:RequiredOrderings}).

\begin{figure}[h]
\sporkCodeRule
\centering
\begin{sporkfigure}
  % Body of the Required Orderings figure (26.tex frame RequiredOrderings, paper RequiredOrderings.tex):
% raw TikZ, \input inside a sporkfigure. Mutates in a row with the read between them;
% reveals match the cases: 2 mutate chain, 3 mutates -> read and read -> mutates, 4 squeeze note,
% 5 hazard boxes. Node names m0, m1, r, m2, m3 are its interface.
\tikzset{
  req event/.style={rectangle, draw=black!40, thick, align=center, inner sep=0.7mm, minimum width=15mm},
  req arrow/.style={spork arrow, dash pattern=on 2.5pt off 2pt},
  req star/.style={font=\fontsize{12}{14}\selectfont, inner sep=0.3mm},
}
\node (m0) [req event] {\sporkMutateBullet\\$\sporkx[\sporkz^*]$};
\node (m1) [req event, right=7mm of m0] {\sporkMutateBullet\\$\sporkx[\sporkz^*]$};
\node (r) [req event, anchor=north west] at ([xshift=14mm, yshift=-3.5mm]m1.east)
  {\sporkReadBullet\\$\sporkx[\sporkz^*]$};
\node (m2) [req event, anchor=west] at ([xshift=14mm, yshift=3.5mm]r.north east)
  {\sporkMutateBullet\\$\sporkx[\sporkz^*]$};
\node (m3) [req event, right=7mm of m2] {\sporkMutateBullet\\$\sporkx[\sporkz^*]$};
\coordinate (rHi) at ([yshift=1.2mm]r.west);
\coordinate (rLo) at ([yshift=-1.2mm]r.west);
\coordinate (rHiOut) at ([yshift=1.2mm]r.east);
\coordinate (rLoOut) at ([yshift=-1.2mm]r.east);

% Step 2: mutates totally ordered.
\begin{scope}[spork reveal=<2->, draw=sporkBlueFg, text=sporkBlueFg]
  \draw [req arrow] (m0) -- node [req star, above] {$*$} (m1);
  \draw [req arrow] (m1) -- node [req star, above] {$*$} (m2);
  \draw [req arrow] (m2) -- node [req star, above] {$*$} (m3);
  \node [anchor=south] at ($(m0.north)!0.5!(m3.north)$)
    {Totally ordered by \sporkBlue{\sporkMutateBullet{} $\rightarrow^*$ \sporkMutateBullet}; \textbf{overwrites previous}};
\end{scope}

% Step 3: labels under the lower L lines, centered between the outer mutate and the read.
% The outer (longer) arrows' stars sit beside their vertical segments to leave room.
\begin{scope}[spork reveal=<3->, draw=sporkVioletFg, text=sporkVioletFg]
  \draw [req arrow] (m0.south) |- node [req star, left, pos=0.25] {$*$} (rLo);
  \draw [req arrow] (m1.south) |- node [req star, above, pos=0.85] {$*$} (rHi);
  \node (rawLabel) [anchor=north] at ([xshift=-2mm, yshift=-0.8mm]$(m0.west |- rLo)!0.5!(rLo)$)
    {Ordered by \sporkViolet{\sporkMutateBullet{} $\rightarrow^*$ \sporkReadBullet}};
\end{scope}
\begin{scope}[spork reveal=<3->, draw=sporkVioletFg, text=sporkVioletFg]
  \draw [req arrow] (rHiOut) -| node [req star, above, pos=0.15] {$*$} (m2.south);
  \draw [req arrow] (rLoOut) -| node [req star, right, pos=0.75] {$*$} (m3.south);
  \node (warLabel) [anchor=north] at ([xshift=2mm, yshift=-0.8mm]$(rLoOut)!0.5!(m3.east |- rLoOut)$)
    {Ordered by \sporkViolet{\sporkReadBullet{} $\rightarrow^*$ \sporkMutateBullet}};
  % Slide show is where to put these handwavey intuition. Below the labels.
\end{scope}
\begin{scope}[spork reveal=<4->]
  \node (squeezeNote) [annotation, text=black, below=4mm of r]
    {Squeezed between back-to-back $\sporkx[\sporkz^*]$ mutates};
\end{scope}
\draw [annotation line, spork reveal=<4->] (squeezeNote.north) -- (r.south);

% Step 5: the classic hazard each case prevents, in yellow boxes; explained orally.
% WAW sits above its edge, left of the star; RAW / WAR sit under their "Ordered by" labels,
% beside the squeeze note, so no box hides a dashed edge or star.
\begin{scope}[spork reveal=<5->]
  \tikzset{hazard/.style={spork yellow style, rectangle, thick, inner sep=0.8mm, font=\bfseries}}
  \node [hazard, anchor=south] at ([xshift=-12mm, yshift=0.8mm]$(m1.east)!0.5!(m2.west)$) {WAW};
  \node [hazard, anchor=north west] at ([yshift=-0.5mm]rawLabel.south west) {RAW};
  \node [hazard, anchor=north east] at ([yshift=-0.5mm]warLabel.south east) {WAR};
\end{scope}

\end{sporkfigure}
\caption{Required agreement of trace total order (left to right) and parallel partial order ($\dashrightarrow^*$) for mutate/mutate and mutate/read pairs of the same $\sporkx[\sporkz^*]$.}
\label{fig:RequiredOrderings}
\sporkCodeRule
\end{figure}


What remains is to summarize the contents of a trace event and what parallel partial order is.
To ease analysis,  % Kind of not true; it's actually strictly necessary to avoid transitivity through instructions.
we split array accesses into ``precondition'' and ``postcondition'' events $\sporkTrPreAccess$ and $\sporkTrPostAccess$.
These record the accessed array element $\sporkx[\sporkz^*]$,
a read-or-mutate tag $\sporkmut$,
and other state described in Section~\ref{sec:TraceChecks}.
Statements and ``read only'' expressions are carefully stepped to correctly append to the trace
(\sporkTrRead, Figure~\ref{fig:SyncTraceSemantics}).
This is done with the checked trace concatenation operator (Section~\ref{sec:TraceConcatChk}),
which checks, upon appending preconditions events $\sporkTrPreAccess$,
that correct parallel partial ordering was established with previous $\sporkTrPostAccess$ events.
Synchronization events are recorded as $\sporkTrFence$ events (non-split barriers)
or pairs of $\sporkTrArrive,\sporkTrAwait$ events (split barriers).

Totally-trace-ordered events $\sporktr_1,\sporktr_2$ are additionally partially-parallel-ordered when they
\begin{enumerate}
  \item share a common thread and are implicitly ``in-order'', hence seqenced by the thread, or
  \item share a common thread and correspond to an async instruction \& specialized barrier, or
  \item are ordered by usage of a common split barrier (e.g. same mbarrier)
\end{enumerate}
To implement the first two cases,
we give each trace event inbound and outbound \emph{observer sets},
where observers $\sporkob: \sporkOb$ consist of a pair of thread ID $\sporktid$  % \emph{thread ID} is super ugly
and \emph{qualitative timeline} $\sporkqtl$.
The latter distinguishes effects of asynchronous instructions, and their specialized barriers,
from ordinary implicitly-in-order instructions (Section~\ref{sec:Qtl}).
Parallel partial ordering between two totally-trace-ordered events $\sporktr_1,\sporktr_2$ is  % not"'are"; "ordering" is singular.
decided based on overlaps between $\sporktr_1$'s outbound observer set and $\sporktr_2$'s inbound observer set (Section~\ref{sec:TraceEdges}).
The added qualitative timeline makes the overlap test stricter and filters out intra-thread orderings not guaranteed by hardware semantics.
For the split barrier case, we define inbound and outbound \emph{pending await sets}, used similarly.



